\section{Related Work}
\label{sec:related}

\paraf{LLM training.}
A lot of efforts have been put to the training of pre-trained LLMs, including
proprietary ones such as GPT-3~\cite{brown2020language},
GPT-4~\cite{openai2023gpt4}, GShard~\cite{lepikhin2020gshard},
PaLM~\cite{chowdhery2022palm}, and many others~\cite{askell2021general,
wei2022finetuned, smith2022using, hoffmann2022training, su2023welm}, as well as
open-source alternatives like OPT~\cite{zhang2022opt},
BLOOM~\cite{scao2022bloom}, Llama~\cite{touvron2023llama1},
Llama-2~\cite{touvron2023llama}. Existing technical reports in the field
predominantly focus on model performance comparisons, leaving out the specific
details of the system infrastructure that makes such training possible.
This paper fills this gap by sharing our experience of end-to-end LLM
pre-training at the scale of over 10,000 GPUs from a systems perspective.

After pre-training, pre-trained base models can be further fine-tuned to adapt
to downstream tasks better. This has led to the emergence of a range of dialogue
models~\cite{alpaca, vicuna2023, koala_blogpost_2023, BELLE} exemplified by
ChatGPT. However, it is worth noting that the computational power and data
requirements for fine-tuning are substantially lower than those for
pre-training. With the application of optimization techniques such as
quantization~\cite{li2020train, xiao2022smoothquant, frantar2022gptq,
dettmers2022llm} and low-rank adaptation~\cite{hu2021lora}, fine-tuning can be
efficiently accomplished with limited resources. 

\parabf{LLM optimizations.}
In addition to the techniques mentioned previously in the paper, there exists a lot of other works targeted at improving the efficiency of LLMs. Sparse or linear attentions~\cite{child2019generating, katharopoulos2020transformers, zhu2021longshort} are proposed to make the memory consumption scales approximately linearly. Several studies
aim to design new architectures rather than conventional
transformer architectures to address the efficiency issue, such as
RWKV~\cite{peng2023rwkv} and RetNet~\cite{sun2023retentive}. 
Many recent studies have been devoted to developing communication acceleration techniques for LLMs. Some works reduce communication traffic using gradient compression~\cite{alistarh2017qsgd} or mixed-precision training~\cite{DBLP:conf/iclr/MicikeviciusNAD18}, while others schedule communication to overlap it with computation. 
Many popular ML frameworks, such as TensorFlow~\cite{abadi2016tensorflow} and PyTorch~\cite{paszke2019pytorch}, enable overlapping communication with backward propagation by default.
Recent works~\cite{jayarajan2019priority, hashemi2019tictac, peng2019generic, bao2020preemptive} further overlap gradient synchronization with forward computation via tensor partitioning, at the cost of extra overhead. Some works~\cite{li2018pipe, chen2022sapipe} introduce fixed staleness to the training pipeline for full overlapping communication and communication. However, the staleness may degrade the model performance.


\parabf{Diagnosis tools in datacenters.} Many diagnosis tools have been developed
to identify and pinpoint hardware and software problems in datacenters.
Pingmesh~\cite{pingmesh} is an active
probing system based on end hosts. Network wide RTT and packet loss and measured
by sending probing ping packets and doing data analysis. Network-wide SLAs are
provided and network problems including packet-blackhole and packet silent drop
are detected.  EverFlow~\cite{everflow}, LossRadar~\cite{lossradar},
NetBouncer~\cite{netbouncer} exploits the capability of switches to diagnose
detailed network problems like network path failures or specific network port
failures. NetBouncer leverages IP-in-IP tunnel techniques to do path probing.
EverFlow requires mirroring network packets to a centralized server to do
debugging. Hostping~\cite{hostping} is a diagnosis system based on
end hosts that focuses on intra-host bottlenecks. It actively senses complex GPU
server PCIe/NVLINK interconnects and does loopback bandwidth and latency tests.

\parabf{Fault tolerance in large-scale distributed systems.}
Fault tolerance has been a major concern in large-scale distributed systems,
where a wide range of hardware and software failures can occur.
Many fault tolerance techniques have been proposed in the past
that cater the needs of different
systems and deployment scenarios.
Reactive fault
tolerance techniques are used to reduce the impact of failures on a system when
the failures occur. There are many techniques in this category such
as Retry~\cite{haider2011fault}, Replication~\cite{haider2011fault},
Checkpointing~\cite{peng2018optimus} and Message
Logging~\cite{tanenbaum2007distributed}.
These techniques incur some system overhead to recover from failures.
Proactive fault tolerance techniques
keep healthy components in place as backups of the faulty
components, obviating the need of recovery from faults and errors, e.g., preemptive
migration~\cite{chakravorty2006proactive, chakravorty2005proactive,
chen2020elastic} and load balancing~\cite{behera2014performance}. However, these
approaches often assume that failures are predictable,
while it is challenging for real
large-scale distributed systems to predict the failures due to the complexity
of the systems.

